\documentclass[12pt]{jlreq}
\usepackage{amsmath, amssymb}

\begin{document}

\newcommand{\C}{\mathbb{C}}
\newcommand{\R}{\mathbb{R}}
\newcommand{\Z}{\mathbb{Z}}
\newcommand{\Tr}{\operatorname{Tr}}
\newcommand{\Impart}{\operatorname{Im}}
\newcommand{\AF}{\operatorname{AF}}
\newcommand{\AX}{\operatorname{AX}}
\newcommand{\GL}{\operatorname{GL}}

零行列ではない \(2\) 次の正方行列全体 \(X\) を考える。
\[
  X=\left\{
    A=\begin{pmatrix}a&b\\ c&d\end{pmatrix}
    \,\middle|\,
    a,b,c,d\in\R,\ 
    A\ne \begin{pmatrix}0&0\\0&0\end{pmatrix}
  \right\}.
\]
\(X\) を同値関係 \(A\sim \lambda A\ (\lambda>0)\) で割った商空間を \(P\) とおき，射影を
\(\pi:X\to P\) とおく。
\[
  N=\{A\in X\mid \det A=0\}
\]
に対して \(M=\pi(N)\) とおく。\(X\) にはユークリッド空間の部分空間としての位相をいれ，
\(P\) にはその商位相，\(M\) には \(P\) の部分空間としての位相をいれる。

\begin{enumerate}
  \item 次の写像
  \[
    M\xrightarrow{\iota}P\xleftarrow{\pi}X
  \]
  が共に \(C^\infty\) 写像となるように，\(P\) および \(M\) に \(C^\infty\) 多様体の構造が入ることを示せ。
  ここで \(\iota\) は包含写像である。
  \item \(M\) 上の任意の \(C^\infty\) 関数は，\(P\) 上の \(C^\infty\) 関数に拡張できることを示せ。
  \item \(M\) 上の任意の \(C^\infty\) 関数 \(h\) に対し，(2) により \(P\) 上に拡張した \(C^\infty\) 関数の \(1\) つを \(H\) とする。
  \(H\) の \(\pi\) による引き戻しを \(F\) とおく。\(X\) 上の関数
  \[
    \left(\frac{\partial^2}{\partial a\,\partial d}
      -\frac{\partial^2}{\partial b\,\partial c}\right)F
  \]
  の \(N\) への制限は，拡張 \(H\) のとり方によらず \(h\) のみで定まることを示せ。
\end{enumerate}

\end{document}
